\section{Introduction}

Across 943 held-out BioASQ test questions, a strong Hybrid RRF retriever fails to place any gold evidence passage in the top 10 for over 6\% of queries where that evidence exists somewhere in the top 100. These are not retrieval failures---gold passages are in the candidate pool. They are ranking failures: passages that a medically informed reader would recognize as relevant are scattered below the top-10 cutoff by scoring functions that rely on pairwise similarity alone.

Biomedical evidence is organized through higher-order relations---among questions, passages, MeSH concepts, biomedical entities, documents, and knowledge-base relations---that pairwise comparisons systematically miss. A question about thyroid hormone analogs and a passage about DITPA share no surface lexical overlap, yet a clinician connects them instantly through the MeSH hierarchy: both map to thyroid-related tree nodes. Standard lexical and dense retrieval rank by pairwise query-passage similarity and lose this connection.

We develop KCH-MedRank, a knowledge-constrained learning-to-rank framework that embeds these higher-order relations as interpretable ranking features. The method maintains reproducible candidate generation through BM25, dense retrieval, and Reciprocal Rank Fusion~\citep{robertson2009bm25,cormack2009rrf}, then enriches each query's local candidate neighborhood with retrieval scores, biomedical semantic features, MeSH hierarchy-aware features, entity and relation features, and optional hypergraph structural features. A local hypergraph serves as one architectural variant for capturing n-ary concept interactions; a flat knowledge-feature model and a pairwise graph decomposition serve as controlled baselines.

The core contribution is a focused evidence reranking framework with post-hoc feature-attribution analysis. KCH-MedRank integrates these features under a query-group LambdaMART ranker to produce the final evidence ordering, and the reranker's feature importance scores expose which signals---MeSH hierarchy paths, shared entity clusters, diffusion centrality---drive each ranking decision.

On the BioASQ evidence retrieval benchmark, KCH-MedRank significantly improves Recall@10 over a true MedCPT Cross-Encoder baseline on the same enhanced candidate pool ($+0.0157$, $p=0.0076$). A structural ablation reveals that flat biomedical knowledge features already account for most of the gain, and that adding medical knowledge constraints to the local hypergraph is necessary: without them, raw hypergraph topology introduces noise and degrades ranking. With medical constraints, the degradation reverses with statistically significant top-10 gains ($p < 0.01$). The hypergraph provides additional early-rank benefit at $k{=}5$, complementing the flat knowledge signal at deeper ranks. Supplementary diagnostics on PubMedQA and BEIR NFCorpus confirm evidence coverage improvements and reranking robustness.

Our contributions are: (1)~a knowledge-constrained learning-to-rank framework that combines retrieval, biomedical semantic, MeSH hierarchy, entity, relation, and optional hypergraph features for evidence reranking; (2)~a structural ablation framework that isolates the effects of flat biomedical knowledge, graph topology, and medical knowledge constraints, revealing that the latter are necessary for structure to provide useful rather than noisy signal; and (3)~an interpretability analysis tracing 648 rescued gold passages to three knowledge-grounded mechanisms, with the MeSH hierarchy path dominating at 75.9\% of rescues.
